% Generated draft: requires mathematical and source review.
\documentclass[11pt]{article}
\usepackage[margin=25mm]{geometry}
\usepackage{fontspec}
\setmainfont{Latin Modern Roman}
\usepackage{amsmath,amssymb,mathtools,unicode-math}
\setmathfont{Latin Modern Math}
\usepackage{enumitem,needspace,xurl,hyperref}
\setlength{\parindent}{0pt}
\setlength{\parskip}{5pt}
\setlength{\emergencystretch}{3em}
\newcommand{\N}{\mathbb{N}}
\newcommand{\Z}{\mathbb{Z}}
\newcommand{\Q}{\mathbb{Q}}
\newcommand{\R}{\mathbb{R}}
\newcommand{\C}{\mathbb{C}}
\newcommand{\F}{\mathbb{F}}
\newcommand{\A}{\mathbb{A}}
\newcommand{\PP}{\mathbb{P}}
\newcommand{\E}{\mathbb{E}}
\newcommand{\eps}{\varepsilon}
\newcommand{\dd}{\,\mathrm{d}}
\begin{document}
\section*{UnsolvedMath ID: 700611}
\textbf{Source problem code:} catalogue entry
\subsection{Definitions and mathematical statement}
A group \(G\) is \emph{periodic} if every element has finite order. It is \emph{residually finite} if, for every \(g\in G\setminus\{1\}\), there are a finite group \(Q\) and a homomorphism \(\varphi:G\to Q\) with \(\varphi(g)\ne 1\). For a prime \(p\), a \(p\)-subgroup is a subgroup all of whose elements have order a power of \(p\); a Sylow \(p\)-subgroup means a subgroup maximal under inclusion among the \(p\)-subgroups of \(G\).

The supplied question asks whether there exists a group \(G\) satisfying all of the following conditions:
\begin{enumerate}
\item \(G\) is periodic and residually finite.
\item \(G\) is an \(F^{*}\)-group, in the sense referenced by the source as Archive, Problem 7.42.
\item For every prime \(p\), every Sylow \(p\)-subgroup of \(G\) is finite.
\item \(G\) is not binary finite.
\end{enumerate}
The supplied record does not include the referenced definition of an \(F^{*}\)-group or a definition of \emph{binary finite}. These two conditions retain the source terminology and remain unresolved predicates in this draft; a complete formal statement requires their intended definitions. No uniform bound on the orders of the finite Sylow subgroups is stated.
\subsection{Short English statement}
Does there exist a periodic, residually finite \(F^{*}\)-group with every Sylow subgroup finite, but which is not binary finite? The meanings of \(F^{*}\)-group and binary finite require source clarification. The supplied registry labels the problem open; that label is not independently verified here.
\subsection{Sources}
\begin{itemize}
\item Supplied local UnsolvedMath registry record: permanent ID \(700611\), title \emph{Kourovka Notebook Problem 9.78}, exact source code \texttt{catalogue entry}, upstream revision \texttt{d371ee917ac9647fbb06cd6fc529deec6f3dda82}. The supplied status is \emph{open}.
\item \emph{The Kourovka Notebook (Unsolved Problems in Group Theory)}, Problem 9.78, proposed by V. P. Shunkov, ninth issue (1984); the record identifies the twenty-first issue edition (2026), \texttt{arXiv:1401.0300v48}. The supplied collection description names E. I. Khukhro and V. D. Mazurov as editors.
\item The statement refers to Archive, Problem 7.42, for the meaning of \(F^{*}\)-group. The contents of that archive entry were not supplied. The cited sources have not been independently consulted for this draft.
\end{itemize}
\end{document}
